% !TEX program = xelatex
% 编译方式：在本目录执行两遍 xelatex（详见 docs/_manual_common/README.md）。
% 源码依据：src/ev_power_traffic/ 与 include/hacdcpf/ev_power_traffic/ 源码及
%           docs/guides/ev_traffic_scenario_format.md，公式与参数以代码为准。
\documentclass[UTF8,fontset=fandol,a4paper,11pt]{ctexart}
\usepackage{../../_manual_common/hysim_manual}

\title{\textbf{交直流混合高弹性能源电力系统仿真平台}\\[0.5em]
       {\Large 电动汽车-交通-电力耦合技术手册}\\[0.3em]
       {\large 数学模型、接口参数、验证校核与工业级评价}}
\author{HySim-XJTU-HRPES（西安交通大学 HRPES 团队）}
\date{\today}

\begin{document}
\maketitle
\begin{abstract}
\noindent 本手册依据 HySim-XJTU-HRPES 平台电动汽车-交通-电力耦合模块
（\srcpath{src/ev\_power\_traffic/}、\srcpath{include/hacdcpf/ev\_power\_traffic/}）整理，
覆盖 A--H 八个建模层级：元胞传输模型（CTM）/连接传输模型（LTM）交通流、动态用户均衡（DUE）、
联合社会福利 LP/MILP/MPEC、充电站选址定容 MILP 与滚动时域 MPC，命名空间 \texttt{hacdcpf::evpt}。
诚实口径为本手册核心：顶层入口 \srcpath{simulate\_ev\_power\_traffic} 明确为启发式/分解式，
其 \fld{mathematical\_model\_verified}=false；每个选项结构均暴露 \fld{require\_exact\_mathematical\_model}，
在开启时\textbf{拒绝}返回伪装成证明的启发式结果，而非默默降级。全部结论以源码与
\srcpath{docs/guides/ev\_traffic\_scenario\_format.md} 为准。文档版式与《碳流追踪与碳排放核算技术手册》一致。
\end{abstract}

\tableofcontents
\newpage

\section{阅读指南与约定}
\subsection{数据来源}
\begin{itemize}
  \item \srcpath{include/hacdcpf/ev\_power\_traffic/simulation.hpp} —— A--H 各建模层入口；
  \item \srcpath{types.hpp}、\srcpath{options.hpp}、\srcpath{joint\_social\_welfare.hpp}、
        \srcpath{route\_generation.hpp}、\srcpath{ltm\_network.hpp} 及总入口
        \srcpath{ev\_power\_traffic\_simulation.hpp} —— 类型、选项与子模型；
  \item \srcpath{src/ev\_power\_traffic/}（\srcpath{ctm\_propagation.cpp}、\srcpath{ctm\_so\_lp.cpp}、
        \srcpath{ctm\_due\_vi.cpp}、\srcpath{ltm\_propagation.cpp}、\srcpath{ltm\_so\_lp.cpp}、
        \srcpath{ltm\_mpc.cpp}、\srcpath{joint\_optimizer.cpp}、
        \srcpath{joint\_optimizer\_full\_nlp.cpp}、\srcpath{infra\_design\_milp.cpp} 等）—— 实现；
  \item \srcpath{docs/guides/ev\_traffic\_scenario\_format.md} —— 场景格式与口径契约；
  \item \srcpath{tests/test\_ev\_power\_traffic*.cpp} 等 —— 验证依据。
\end{itemize}
命名空间为 \texttt{hacdcpf::evpt}。

\subsection{单位制与索引空间}
时间以小时为单位，CTM 步长 \fld{dt\_ctm\_hr} 受 CFL 条件约束；路段长度 km、
容量 veh/hr、能耗 kWh/(veh$\cdot$km)、功率 MW、电价 \$/MWh。交通图（\texttt{TrafficGraph}）
的元胞/路段索引与电网侧 \texttt{HybridPowerSystem} 的母线索引在 DC-OPF 耦合处按站-母线映射对齐；
路径专属（route-specific）元胞在分流处保持出行（站点选择）身份，索引不跨域裸传。

\subsection{诚实口径}
顶层 \srcpath{simulate\_ev\_power\_traffic} 为启发式/分解式，
\texttt{EVPowerTrafficResult.}\fld{mathematical\_model\_verified}=false，状态标注
“simulation entry includes heuristic or decomposed layers”。DUE 以不动点/变分不等式（VI）
残差判定收敛（\fld{relative\_gap}、\fld{vi\_gap}），非全局求解器证明；充电调度是否贪心由
\fld{charging\_dispatch\_is\_greedy} 显式标注。认证型 MPEC-MILP 模式仅对其文档化的
有限/线性/PWL 模型全局认证；\texttt{FullJoint*NLP} 模式仅携带局部驻点证书，在
\fld{require\_global\_nonlinear\_certificate} 下返回不支持。聚合型 CTM/LTM 连接变量模式为
降维松弛。开启 \fld{require\_exact\_mathematical\_model} 时，任何启发式/近似路径均被拒绝。

\section{功能定位与架构}
\subsection{公共入口族（建模层 A--H）}
\begin{center}\footnotesize
\begin{tabular}{@{}L{6.8cm}L{2.4cm}L{5.8cm}@{}}
\toprule
入口 & 层级 & 说明\\
\midrule
\srcpath{simulate\_ev\_power\_traffic(problem, options)} & A & 顶层启发式/分解式（分配+充电+可选潮流）\\
\srcpath{simulate\_ev\_power\_traffic\_ctm\_due(...)} & B & CTM + 动态用户均衡（MSA）\\
\srcpath{simulate\_ev\_power\_traffic\_ctm\_joint(...)} & C & CTM-DUE $\leftrightarrow$ DC-OPF 的 LMP 价格协调\\
\srcpath{solve\_joint\_optimizer(problem, options)} & D & 单一 LP/MILP（路径+充电/V2G+DC-OPF；认证 MPEC-MILP）\\
\srcpath{solve\_ctm\_so\_lp(...)} / \srcpath{solve\_ltm\_so\_lp(...)} & E & 系统最优 CTM/LTM 线性规划\\
\srcpath{solve\_ctm\_due\_vi(...)} & F & Beckmann Frank-Wolfe 的 DUE-VI\\
\srcpath{solve\_infra\_design\_milp(...)} & G & 充电站选址/插座定容 MILP\\
\srcpath{solve\_ltm\_mpc(...)} & H & LTM 滚动时域 MPC\\
\srcpath{solve\_joint\_social\_welfare(problem, options)} & D & 社会福利联合优化\\
\srcpath{generate\_candidate\_routes(...)} / \srcpath{simulate\_ltm\_network(...)} & — & 候选路径生成 / LTM 网络仿真\\
\bottomrule
\end{tabular}
\end{center}

\subsection{关键数据结构}
\compmeta{EVPowerTrafficProblem}{include/hacdcpf/ev\_power\_traffic/simulation.hpp}
主要字段：\fld{system}（\texttt{HybridPowerSystem}）、\fld{traffic}（\texttt{TrafficGraph}）、
\fld{routes}、\fld{demands}、\fld{station\_prices}、\fld{icv\_demands}、\fld{station\_spillback}。

\compmeta{EVPowerTrafficOptions}{include/hacdcpf/ev\_power\_traffic/options.hpp}
主要字段：\fld{num\_steps}、\fld{assignment\_model}、\fld{allow\_v2g}、
\fld{auto\_generate\_routes}/\fld{k\_shortest\_paths}、\fld{require\_exact\_mathematical\_model}；
子结构 \texttt{CTMOptions}（\fld{dt\_ctm\_hr}、\fld{n\_cells\_per\_link}、
\fld{route\_specific\_cells}、\fld{enable\_multiclass}+\fld{ev\_pce}/\fld{icv\_pce}、
\fld{use\_aggregate\_link\_variables}）与 \texttt{JointOptimizerOptions}
（\texttt{JointOptimizerMode}：\texttt{SocialWelfareMax}、\texttt{IntegerRouteMILP}、
\texttt{CertifiedDynamicMPECMILP}、\texttt{CertifiedFullJointLtmPwlMILP}、
\texttt{FullJointSocial\allowbreak WelfareNLP} 等；\fld{include\_dcopf}、
\fld{allow\_v2g}、\fld{mip\_gap}、
\fld{time\_limit\_sec}、\fld{full\_joint\_ltm\_pwl\_segments}）。

\compmeta{TrafficLink}{include/hacdcpf/ev\_power\_traffic/types.hpp}
主要字段：\fld{length\_km}、\fld{free\_flow\_time\_hr}、\fld{capacity\_veh\_per\_hr}、
\fld{jam\_vehicles}、BPR \fld{alpha}=0.15/\fld{beta}=4、\fld{drive\_energy\_kwh\_per\_veh\_km}。
结果结构 \texttt{CTMDUEResult}（\fld{converged}、\fld{relative\_gap}、\fld{vi\_gap}、
\fld{mathematical\_model\_verified}、\fld{charging\_dispatch\_is\_greedy}、
\fld{flow\_by\_demand\_route}）、\texttt{InfraDesignResult}（\fld{station\_open}、
\fld{station\_plugs}、\fld{mip\_gap}）与 \texttt{LTMMPCResult}（\fld{admission}）。

\section{源码等价的 A--H 执行模型}
\label{sec:evpt-source-equivalent}

本章按公开入口分别记录实现，禁止把不同入口拼成一个不存在的统一优化模型。所有严格模式均保留源码
拒绝语义：\texttt{simulate\_ev\_power\_traffic} 的路线分配、充电调度和电网验证不是一个全局认证
程序，\texttt{require\_exact\_mathematical\_model=true} 时直接返回不支持。依据：
\srcpath{src/ev\_power\_traffic/ev\_power\_traffic\_simulation.cpp:61--95}。

\subsection{A：分解仿真入口}
入口为 \texttt{simulate\_ev\_power\_traffic}。它先生成或读取候选路径，再按
\texttt{assignment\_model} 执行路径分配，生成充电会话并逐时验证电网。最终
\texttt{feasible} 的源码谓词是：系统最优求解请求必须求解成功；若禁止未服务出行，则未服务车辆
不超过 $10^{-7}$；若禁止未服务充电，则未服务电量不超过 $10^{-6}$；每步发电容量必须合格，且
已经运行的潮流必须收敛。依据：
\srcpath{src/ev\_power\_traffic/ev\_power\_traffic\_simulation.cpp:7--41}。

该入口不是单体数学规划，因此本章不为 A 添加联合目标函数。站点负荷写回只执行
\begin{equation}
 p_s^{model}=p_s^{EV},\qquad q_s^{model}=q_s^{EV},\qquad
 utilization_s=\begin{cases}|p_s^{EV}|/P_s^{max},&P_s^{max}>10^{-9},\\0,&\text{否则。}\end{cases}
\end{equation}
依据：\srcpath{src/ev\_power\_traffic/ev\_power\_traffic\_simulation.cpp:46--58}。

\subsection{B/F：CTM 传播与 DUE--VI}
对路段 $a$，若长度和 CTM 步长有效，元胞数为
\begin{equation}
 M_a=\begin{cases}
 \texttt{n\_cells\_per\_link},&\text{该选项}>0,\\
 \max(1,\lfloor L_a/(v_a^f\Delta t_c)\rfloor),&\text{否则，}
 \end{cases}
 \qquad\delta_a=L_a/M_a.
\end{equation}
无效长度或步长时固定 $M_a=1$ 且 $\delta_a=\max(L_a,10^{-9})$。容量非正时取
$10^4$ veh/h。拥堵波速按
\begin{equation}
 k^{jam}=\begin{cases}N^{jam}/\delta,&N^{jam}>10^{-9},\\2q^{max}/v^f,&\text{否则，}\end{cases}
 \quad k^{crit}=q^{max}/v^f,
 \quad w=\begin{cases}q^{max}/(k^{jam}-k^{crit}),&k^{jam}>k^{crit}+10^{-9},\\
 w^{fallback},&\text{否则。}\end{cases}
\end{equation}
依据：\srcpath{src/ev\_power\_traffic/ctm\_propagation.cpp:40--103}。

元胞发送、接收和内部边界流严格为
\begin{align}
 S_{a,m,k}&=\Delta t_c\min(n_{a,m,k}v_a^f/\delta_a,q_{a,k}^{max}),\\
 R_{a,m,k}&=\Delta t_c\min(q_{a,k}^{max},
 w_a\max(0,N_a^{jam}-n_{a,m,k})/\delta_a),\\
 y_{a,m,k}&=\min(S_{a,m,k},R_{a,m+1,k}),\\
 n_{a,m,k+1}&=\max(0,n_{a,m,k}+y_{a,m-1,k}-y_{a,m,k}).
\end{align}
路径专属元胞开启时，跨边界的各路径流按上游旧占有量比例分配，而非独立求最小值。节点合流/分流
还有两阶段预算缩放，不能用上述内部元胞式替代；权威实现为
\srcpath{src/ev\_power\_traffic/ctm\_propagation.cpp:401--752}。内部元胞依据：
\srcpath{src/ev\_power\_traffic/ctm\_propagation.cpp:110--127}、
\srcpath{src/ev\_power\_traffic/ctm\_propagation.cpp:313--385}。

F 入口为 \texttt{solve\_ctm\_due\_vi}。精确模型开关为真时直接拒绝，因为结果只给固定点残差，
不是全局 VI/MPEC 证明。仿真步与 CTM 步的倍率是
\begin{equation}
 R=\max(1,\operatorname{round}(\Delta t_{sim}/\Delta t_c)),\qquad T_c=T_{sim}R.
\end{equation}
初始流在每个需求的出发步集合和可行候选路径间等分。混合车流开启时，送入 CTM 的流为
\begin{equation}
 x^{combined}_{k,r}=PCE_{EV}x^{EV}_{k,r}+PCE_{ICV}x^{ICV}_{k,r}.
\end{equation}
依据：\srcpath{src/ev\_power\_traffic/ctm\_due\_vi.cpp:66--106}、
\srcpath{src/ev\_power\_traffic/ctm\_due\_vi.cpp:166--236}。

每轮全有或全无方向为 $\mathbf d=\mathbf y-\mathbf x$。若
\texttt{line\_search\_steps}>0，代码在 $[0,1]$ 上用黄金分割最小化 CTM 重放得到的
$TSTT=\sum x_{k,r}t_{k,r}$；否则
\begin{equation}
 \alpha_i=\frac1{i+1},\qquad\mathbf x\leftarrow\mathbf x+\alpha_i\mathbf d.
\end{equation}
这不是旧文档所写的 $1/i$。收敛先按 Wardrop 相对间隙小于
\texttt{convergence\_tol} 判定；最终若 VI 证书可用，报告间隙还取既有报告、EV/ICV 间隙与
归一化 VI 间隙的最大值。混合车流时 VI 证书只覆盖 EV 子博弈。依据：
\srcpath{src/ev\_power\_traffic/ctm\_due\_vi.cpp:239--430}。

\subsection{C：价格协调，不是单体最优}
C 每轮先运行 CTM--DUE，再逐时把站点有功以 $p_s^{EV}/1000$ MW 加入母线负荷并求 DC--OPF。
OPF 收敛时，站价更新为
\begin{equation}
 \pi^{new}_{s,k}=\alpha\frac{LMP_{bus(s),k}}{1000}+(1-\alpha)\pi^{old}_{s,k},
 \qquad\alpha=\operatorname{clamp}(\texttt{price\_update\_step},0,1).
\end{equation}
OPF 失败的步保留旧价。收敛仅指
\begin{equation}
 \max_{s,k}|\pi^{new}_{s,k}-\pi^{old}_{s,k}|<
 \texttt{price\_convergence\_tol};
\end{equation}
它不证明社会福利最优。报告福利按代码后处理为
$W=B_{EV}-C_{gen}-VOT\cdot TSTT$。禁用 DC--OPF 时固定价格且只运行一轮，结果被置为收敛；
此时报告生成成本为零而不是静态估计值。依据：
\srcpath{src/ev\_power\_traffic/ctm\_joint\_welfare.cpp:237--355}。

\subsection{D：按模式分派的联合程序}
入口 \texttt{solve\_joint\_optimizer} 根据 \texttt{JointOptimizerMode} 分派到不同模型：基础
LP/MILP、有限路线 BPR NLP/MPEC、认证动态 MILP 或全联合 LTM--PWL MILP。它们的变量、约束和
证书不同，不能共享一条目标式。基础模型的权威装配为
\srcpath{src/ev\_power\_traffic/joint\_optimizer.cpp:1652--2644}；非线性模式为
\srcpath{src/ev\_power\_traffic/joint\_optimizer\_full\_nlp.cpp:1760--1915}。

结果字段已区分 \texttt{proven\_optimal}、\texttt{local\_optimum\_certificate}、
\texttt{global\_optimum\_certificate}、\texttt{travel\_times\_are\_endogenous}、
\texttt{wardrop\_complementarity\_enforced} 和 \texttt{dcopf\_coupling\_enforced}。
文档不得把局部 KKT、PWL 近似或一次 CTM 预运行的外生阻抗改写为全局联合最优。

\subsection{E 与 E--LTM}
\texttt{solve\_ctm\_so\_lp} 和 \texttt{solve\_ltm\_so\_lp} 是不同的系统最优 LP 入口，
分别在 \srcpath{src/ev\_power\_traffic/ctm\_so\_lp.cpp:15} 与
\srcpath{src/ev\_power\_traffic/ltm\_so\_lp.cpp:39} 装配。前者使用时间展开 CTM 节点/路径连续性；后者
使用累积计数 LTM。由于两份装配的行结构不同，本章不使用 Beckmann 积分式替代其中任一实现。

\subsection{G：约化选址定容 MILP}
严格模型开关为真时 G 直接拒绝，状态说明该 MILP 使用目标/动态近似。实际站点变量为二进制
$z_s$ 和整数 $B_s\in[0,B^{max}]$，并加入
\begin{equation}
 B_s-B^{max}z_s\le0.
\end{equation}
每个有效需求建立
\begin{equation}
 \sum_{r\in R_d}h_{d,r}+u_d=D_d,
\end{equation}
其中路径流的初始目标系数为
$VOT\,t^{ff}_{d,r}\cdot\texttt{operating\_cost\_scale}$，未服务变量系数为
$\max(0,VOT\,T)$。站点逐步功率和插头约束实际为
\begin{align}
 \sum_{d,r\ni s}p^{default}h_{d,r}-P_s^{default}z_s&\le0,\\
 \sum_{d,r\ni s}h_{d,r}-B_s&\le0.
\end{align}
依据：\srcpath{src/ev\_power\_traffic/infra\_design\_milp.cpp:35--201}。

若启用道路容量，G 并未精确积分 BPR；它创建 $f_a,e_a$，使用固定
$\alpha=0.15,\beta=4$ 和 $\alpha\beta=0.6$：
\begin{align}
 f_a&=\sum_{d,r:a\in r}h_{d,r},\\
 f_a-e_a&\le cap_a,\qquad f_a\le2cap_a,
\end{align}
并给 $f_a$、$e_a$ 分别赋 $VOT\,t_a^0\,scale$ 与其 $0.6$ 倍的线性系数。启用 LTM 动态行时，
这些静态路径/BPR 代理系数被清零，改由累计计数 VHT 系数承载。依据：
\srcpath{src/ev\_power\_traffic/infra\_design\_milp.cpp:203--288}、
\srcpath{src/ev\_power\_traffic/infra\_design\_milp.cpp:290--370}。

\subsection{H：滚动 LTM LP}
H 每个仿真步只执行当前滚动窗 LP 的第一个控制量。对站点，队列、能量和 V2G 行为为
\begin{align}
 Q_{j+1}-Q_j-b_j+c_j&=0,\qquad c_j\le Q_j,\\
 SOC_{j+1}-SOC_j-Eb_j+Ec_j+\Delta t\,v_j^{2g}&=0,\\
 v_j^{2g}-0.8SOC_j&\le0.
\end{align}
目标系数只对 $Q_{j+1}$ 加 \texttt{weight\_queue}，对各跟踪路段的
$N^{in}_{j+1}-N^{out}_{j+1}$ 加 \texttt{weight\_travel\_time}$\Delta t$，对 V2G 加
$-\texttt{weight\_v2g}$。LTM 行包括累计计数单调、发送/接收时滞、逐步容量，以及
$b_j\le N^{out}_{access,j+1}-N^{out}_{access,j}$。依据：
\srcpath{src/ev\_power\_traffic/ltm\_mpc.cpp:216--419}。

H 优先使用 HiGHS LP，失败后用本地单纯形；只有求解成功才推进持久队列、SOC 和累计计数。
\texttt{proven\_optimal} 还要求原始与整数残差均不超过 $10^{-7}$。依据：
\srcpath{src/ev\_power\_traffic/ltm\_mpc.cpp:423--514}。

\subsection{覆盖边界}
本章逐式覆盖 CTM 内部元胞、F 迭代、C 价格固定点、G 核心约束和 H 滚动行。D、E、E--LTM
规模较大且按模式分支，本章只定位其权威装配和证书语义；没有使用通用 CTM、Beckmann 或 MPEC
方程冒充未逐行转写的源码。


\section{接口与参数}
\begin{paramtable}
\fld{num\_steps} & $T$ & 步 & 整数 & — & 时域离散步数\\
\fld{dt\_ctm\_hr} & $\Delta t$ & h & CFL 界内 & — & CTM 时间步（受 CFL 约束）\\
\fld{n\_cells\_per\_link} & — & — & 整数 & — & 每条路段的元胞数\\
\fld{ev\_pce} & — & — & 正数 & $1$ & EV 折算当量系数（PCE）\\
\fld{icv\_pce} & — & — & 正数 & $1$ & 燃油车折算当量系数\\
\fld{alpha} & $\alpha$ & — & — & $0.15$ & BPR 阻抗系数\\
\fld{beta} & $\beta$ & — & — & $4$ & BPR 阻抗指数\\
\fld{k\_shortest\_paths} & $K$ & — & 整数 & — & 候选路径 K-最短路数\\
\fld{allow\_v2g}（\texttt{EVPowerTrafficOptions}） & — & — & 布尔 & true & 是否允许 V2G 反向馈电（\srcpath{options.hpp:32}）\\
\fld{allow\_v2g}（\texttt{JointOptimizerOptions}） & — & — & 布尔 & false & 是否允许 V2G 反向馈电（Formulation D LP 放电列；false 时放电列上界为零，\srcpath{options.hpp:285}）\\
\fld{mip\_gap} & — & — & $\ge0$ & — & MILP 相对间隙容差\\
\fld{time\_limit\_sec} & — & s & 正数 & — & 求解时限\\
\fld{full\_joint\_ltm\_pwl\_segments} & — & — & 整数 & — & LTM 全联合 PWL 分段数\\
\fld{drive\_energy\_kwh\_per\_veh\_km} & — & kWh/(veh$\cdot$km) & 正数 & — & 行驶单位能耗\\
\fld{require\_exact\_mathematical\_model} & — & — & 布尔 & false & 拒绝启发式/近似的严格开关\\
\end{paramtable}
\fld{require\_exact\_mathematical\_model} 为跨全部选项结构的统一门控：开启后遇启发式/降维/
局部证书路径即返回拒绝，绝不静默降级。

\section{验证与边界局限}
\subsection{验证与校核}
注册测试覆盖各建模层与规模：\srcpath{test\_ev\_power\_traffic}、
\srcpath{test\_ev\_power\_traffic\_ctm\_due}、\srcpath{test\_ev\_power\_traffic\_ctm\_joint}、
\srcpath{test\_ev\_power\_traffic\_joint\_opt}、\srcpath{test\_ev\_power\_traffic\_joint\_opt\_d}、
\srcpath{test\_ev\_power\_traffic\_formulations\_e\_f\_g\_h}、
\srcpath{test\_ev\_power\_traffic\_phase1}、\srcpath{test\_ev\_power\_traffic\_largescale}、
\srcpath{test\_evpt\_fidelity\_bound}、\srcpath{test\_evpt\_approximation\_error}、
\srcpath{test\_ltm\_traffic}、\srcpath{test\_ctm\_ltm\_comparison}、
\srcpath{test\_tntp\_parser}、\srcpath{test\_sioux\_falls\_ieee123\_case} 等。

\subsection{边界与局限（如实标注）}
\begin{itemize}
  \item 顶层 \srcpath{simulate\_ev\_power\_traffic} 明确为启发式/分解式，
        \fld{mathematical\_model\_verified}=false（状态注明含启发式/分解层）；
  \item 每个选项结构均暴露 \fld{require\_exact\_mathematical\_model}，开启时\textbf{拒绝}
        返回伪装成证明的启发式结果，而非默默降级；
  \item DUE 以不动点/VI 残差判定，非全局求解器证明；充电调度贪心与否由
        \fld{charging\_dispatch\_is\_greedy} 标注；
  \item 认证 MPEC-MILP 仅对其文档化的有限/线性/PWL 模型全局认证；
        \texttt{FullJoint*NLP} 仅携局部驻点证书，\fld{require\_global\_nonlinear\_certificate}
        下返回不支持；
  \item 聚合型 CTM/LTM 连接变量模式为降维松弛；混合车队 VI 默认仅覆盖 EV，
        除非 \fld{include\_icv\_in\_due}。
\end{itemize}
\section{跨模块工业级技术评价基线}

本章规定本手册所述模块进入工程算例、批量研究或生产服务前必须共同满足的评价基线。
具体模型公式、变量、约束和专用算法以正文相应章节为准；本章不扩大源码已经实现的范围。

\subsection{输入、输出、边界条件与数据结构审计}
输入审计应逐项检查有限性、单位、上下界、枚举值和引用完整性。系统对象必须先按所需层级完成
校验；AC 与 DC 母线采用域限定映射，组件稳定 \texttt{index}、向量位置和临时图索引不得混用。
输出审计必须记录单位、索引空间、模型范围、收敛或可行状态、后端、容差、迭代/节点计数和告警。
空系统、空时域、孤岛、重复 ID、零容量、退化约束与不可用可选后端均属于边界测试集。

\subsection{工程假设与数值稳定性分析}
模型使用的平衡、线性化、稳态、独立性、概率分布、时间离散或网络等值假设必须与应用场景匹配。
数值评价至少包含变量缩放、绝对/相对残差、可行性违反、条件数或枢轴质量、步长/阻尼、迭代停滞
和随机种子复现性。若采用正则化、平滑、回退、松弛或近似，应同时报告触发条件、参数值、对主指标
的敏感性以及是否改变模型口径；不得用“已返回结果”替代收敛或有效性证明。

\subsection{复杂度与资源分析}
令 $n_x$、$n_c$、$n_t$、$n_s$、$|V|$、$|E|$ 分别表示变量、约束、时段、场景、节点和边数。
报告应分别给出建模、求解、后处理的时间与内存。图遍历通常为 $O(|V|+|E|)$；逐时段/逐场景
流程至少线性增长为 $O(n_t n_s)$ 次子问题调用；稠密线性代数最坏可达 $O(n_x^3)$，稀疏方法则应
报告结构非零数、填充率和因子分解峰值内存；MILP/NLP 不给出虚假的多项式最坏界，而报告节点数、
间隙、时限和实测扩展曲线。复杂度结论必须基于固定硬件、固定后端和可复现算例族。

\subsection{异常工况处理}
非法输入应在进入求解器前返回结构化错误；不可行、无界、不收敛、时限、内存不足和后端缺失必须
相互区分。部分结果只有在对应 \fld{ValidityFlags}、\fld{model\_scope}、
\fld{model\_limitations} 或等价字段允许时才可使用。异常路径不得静默丢弃 AC/DC 资产、制造平衡
节点、把 NaN 改写为零或把近似解标为全局最优；系统原始对象应保持可恢复和可重试。

\subsection{与其他模块的接口关系}
接口链路遵循“工程语义模型 $\to$ 校验 $\to$ 投影/装配 $\to$ 求解或分析 $\to$ 结果回投影”。
跨模块传递时保留稳定 ID、单位、符号、时间基准和场景身份；消费潮流、OPF、动态或可靠性结果前，
先检查其收敛、覆盖范围和有效性标志。HTTP/JSON/Python 层只做语义保持适配，不重新解释求解结果。

\subsection{测试与验证建议}
最低验证矩阵包括：手算小例、标准算例、边界/异常输入、随机性质测试、算法或后端交叉校核、
序列化往返、稳定 ID 回投影、AC/DC 同号母线、重复运行确定性和规模扩展测试。数值算法还应加入
残差独立重算、雅可比有限差分或自动微分校核；近似模型应与高保真模型比较并给出误差分布，而非
只报告单个平均值。回归报告必须列明构建配置、算例版本、容差、通过/失败/跳过数及未验证范围。

\subsection{工业级评估指标}
\begin{itemize}
  \item \textbf{正确性}：守恒残差、约束最大违反、解析/基准误差、结果回投影一致性；
  \item \textbf{鲁棒性}：收敛率、不可行识别率、异常分类准确率、扰动敏感性与随机复现率；
  \item \textbf{性能}：端到端时延、吞吐量、峰值内存、并行效率与规模增长斜率；
  \item \textbf{可追溯性}：输入模式版本、稳定 ID 覆盖率、模型范围/限制字段完整率、告警可定位率；
  \item \textbf{工程有效性}：与标准、外部引擎或现场数据的偏差，以及误判/漏判对业务指标的影响。
\end{itemize}
指标应预先给出验收阈值和失败处置；未达到阈值时只能标记为实验性或受限适用。

\subsection{适用范围、局限性与后续改进}
本手册只评价当前源码覆盖的模型、后端和数据结构，不对未建模物理过程、未安装求解器或未执行的
外部验证作保证。后续改进应按“缺口 $\to$ 可量化指标 $\to$ 源码/测试变更 $\to$ 独立验证”闭环，
优先消除会造成错误结论的语义缺口，其次提升数值鲁棒性与性能，最后扩展模型范围。


\end{document}
